\section{historical notes}

\label{nota:Peano}%
\index{Peano!Giuseppe}%
\emph{Giuseppe Peano} (1858--1932), a mathematician from Turin, contributed to laying the foundations of mathematical logic. The notation $\exists$ for the existential quantifier is attributed to him. Peano's original definition took $1$ as the first natural number, but in modern mathematics, it is more convenient to include $0$ among the natural numbers, just as the empty set is considered among sets.

\label{nota:Galileo}%
\index{Galileo!Galilei}%
\emph{Galileo Galilei} (1564--1642) observed that perfect squares: $1,4,9,16,\dots$ are, on one hand, a small part of all natural numbers (these numbers become increasingly spaced apart), but on the other hand, they are as numerous as the natural numbers because the correspondence $n\mapsto n^2$ is bijective.

\label{nota:Cantor}%
\index{Cantor!Georg}%
\label{nota:Russell}%
\index{Russell!Bertrand}%
\label{nota:Frege}%
\index{Frege!Gottlob}%
Set theory was introduced by \emph{Georg Cantor} (1845--1918) without a true
logical formalization (today we would call it \emph{naive set theory}).
\emph{Gottlob Frege} (1848--1925) was the first mathematician who attempted to
formalize Cantor's set theory. In 1902, \emph{Bertrand Russell}, having read
Frege's work, sent him a letter stating the paradox he had discovered: ``I am in
complete agreement with you on all essential points, particularly when you
reject any psychological element, giving great value to the ideography for the
foundation of mathematics and formal logic [...] there is only one point where I
have encountered a difficulty [...]''. Frege's response (June 22, 1902) is
disheartening: ``Your discovery of the contradiction has caused me great
surprise and, I would say, dismay, because it has shaken the foundations upon
which I intended to build arithmetic.''

\label{nota:Euclide}%
\label{nota:Hilbert}%
\index{Euclide}%
\index{axioms!of Euclid}%
\index{Hilbert!David}%
\index{Euclidean geometry}%
\emph{Euclid} (circa 300 B.C.) introduced the axiomatic method in geometry. His treatise ``The Elements'' is considered the book that has had the greatest impact in the history of mathematics. The axioms introduced by Euclid correspond to geometric constructions made with a straightedge and compass. If we consider the set of points that can be constructed from a finite set of given points, we obtain a dense set that is not complete. It is well known, in fact, that with a straightedge and compass, it is not possible to perform the \emph{trisection of an angle} (i.e., divide a generic angle into three equal parts), nor the \emph{duplication of the cube} (i.e., construct the side of a cube whose volume is double that of a given cube), nor the \emph{squaring of the circle} (i.e., construct the side of a square with the same area as a given circle). The first two problems require the construction of cube roots, while the third problem requires the construction of $\sqrt\pi$. Through Galois theory, it has been demonstrated that constructions with a straightedge and compass can only construct those numbers that can be expressed from integers using, in addition to the four operations, the extraction of square roots. It is easy to construct, for example, the diagonal of a square with a given side: this corresponds to the construction of $\sqrt 2$. It is obvious, however, that transcendental numbers (like $\sqrt\pi$) cannot be constructed, since the permissible operations do not allow us to leave the set of algebraic numbers. The construction of a regular polygon inscribed in a circle of a given radius is equivalent to the construction of the $n$-th complex roots of unity and is thus related to the constructibility of particular algebraic numbers. Only thanks to the contribution of Gauss (1777-1855) was there an advancement over Euclid's knowledge regarding which regular polygons are constructible. Surprisingly, while regular polygons with $7$, $9$, $11$, and $13$ sides are not constructible with a straightedge and compass, the regular polygon with $17$ sides is constructible (Gauss-Wentzel).

This entire discussion should make it clear that it is not obvious how one might
define geometric space rigorously so that it is complete. Indeed, only in 1900
did \emph{David Hilbert} (1862--1943) provide a rigorous axiomatic definition of
Euclidean geometry. The property of completeness is captured by Hilbert through
a maximality axiom: Euclidean space is a set of points that satisfy certain
axioms (the usual postulates of Euclid rearranged by Hilbert) and is also the
largest set of points with such properties. Thus, it must contain the limits of
Cauchy sequences (or the separation points of separated linear sets, if we think
of the continuity of the order instead of completeness) because otherwise, such
points could be added to the space without violating the other axioms but
thereby violating maximality.

\label{nota:Dedekind}
\index{Dedekind!Richard}%
Richard Dedekind (1831--1916), to explain the motivations that led him to find a
rigorous definition of real numbers, writes: ``It is often claimed that
differential calculus deals with continuous quantities, and yet, an explanation
of this continuity is never given; even the most rigorous expositions of
differential calculus do not base their proofs on continuity but, more or less
consciously, appeal either to geometric notions or those suggested by geometry,
or depend on theorems that have not been proven in a purely arithmetic manner.''